% Nouveau EGA V (cote 25): archive cote 25, site page 20 (PDF page 21).
% typescript. Draft compared: none.
% AI-transcribed by Claude Opus 5.5 (claude-opus-5-5, adaptive thinking, effort high), Message Batches API, batch ega25b (msgbatch_01A7n5YuXubUPNK1zsepi8RG), 2026-11-13; not yet checked by eye.
\scanpage{25/20}
\mspage{\hand{17}}

\uline{Proposition} 4.3.    Soit P l'une des propriétés  suivantes (pour un
préschéma algébrique sur un corps K):

(i)  Z \var{\xout{satisfait}}{}{typewriter strike-out}{ega25_25_020-v0001} est lisse sur K.

(ii)  Z satisfait la propriété (R$_k$) géométrique\var{\xout{x}}{}{typewriter strike-out}{ega25_25_020-v0002} sur K.

(iii)  Z est séparable sur K.

(iv)  Z est géom normal sur K.

\var{\margin{$\uparrow$}}{\placedmargin{$\uparrow$}}{marginal note placed}{ega25_25_020-v0003}(\var{\over{v}{\unsure{vi|v}}}{\unsure{vi|v}}{writing over}{ega25_25_020-v0004})  Z  est géom. irr.  sur K.

% NOTE: The following overwritten label is circled in pen.
(\var{\over{vi}{\unsure{v|vi}}}{\unsure{v|vi}}{writing over}{ega25_25_020-v0005})  Z est géom. intègre  sur K.

Soit E l'ensemble des \hand{$\xi \in$} $\overset{\hand{\vee}}{P}$  tels que $Y_{\hand{\xi}}$ satisfasse à P. Alors

a)  E est une partie constructible de $\overset{\hand{\vee}}{P}$ .

b)  Supp\var{\over{x}{o}}{o}{writing over}{ega25_25_020-v0006}sons \var{\xout{dx}}{}{typewriter strike-out}{ega25_25_020-v0007} S=Spec(k), k un corps\var{\xout{x}}{}{typewriter strike-out}{ega25_25_020-v0008}, et supposons \var{\inserted[typed]{dans les cas (i) à (v)}}{dans les cas (i) à (v)}{insertion}{ega25_25_020-v0009} que k soit de
car zéro  ou que f:X$\hand{\rightarrow}$P  soit \var{\deleted{une immersion}}{}{deletion}{ega25_25_020-v0010}\var{\inserted{\unsure{un ramifié|non ramifié}}}{\unsure{un ramifié|non ramifié}}{insertion}{ega25_25_020-v0011}. Supposons enfin dans les
cas (v ) et (vi)  que dim $\hand{\overline{\textcolor{black}{f(X)}}}$ \var{\over{$\hand{>}$}{$\geqslant$}}{$\geqslant$}{writing over}{ega25_25_020-v0012} 2.  Supposons que X satisfasse à P,
alors E contient un ouvert \var{\xout{non vide}}{}{typewriter strike-out}{ega25_25_020-v0013} dense de $\overset{\hand{\vee}}{P}$ .

Démonstration identique à celle de 4.2.   Remarquons que dans
les cas (i) à (\unsure{\var{\deleted{i}}{}{deletion}{ega25_25_020-v0014}v}), il suffit de supposer même (dans b) que  f(Z) est
fini, où Z est l'ensemble des points de X où \var{\xout{lx}}{}{typewriter strike-out}{ega25_25_020-v0015} P est en défaut.

\uline{Corollaire} 4.4.    Considérons aussi la propriété:

(\var{\xout{x}}{}{typewriter strike-out}{ega25_25_020-v0016}$C_m^{\var{\xout{\unsure{0|o}}}{}{typewriter strike-out}{ega25_25_020-v0017}}$)  \var{\xout{Zxestxx"géométriquementxdexxxxxxx}}{}{typewriter strike-out}{ega25_25_020-v0018} " $\overline{Z}$ ne peut être discon-%
necté par une partie fermée de dimension  \var{\xout{m}}{}{typewriter strike-out}{ega25_25_020-v0019}$\hand{\leqslant}$ m" (où $\overline{Z}$ est $Z\otimes_K\overline{K}$ , $\overline{K}$
la clôture  algébrique de K).

Soit E l'ensemble des \hand{$\xi \in$} $\overset{\hand{\vee}}{P}$  tels que  $Y_{\hand{\xi}}$ sur k(\hand{$\xi$}) satisfasse à $C_m$.
Alors:

a)  E est constructible.

b)  \var{\xout{Si}}{}{typewriter strike-out}{ega25_25_020-v0020} \var{\inserted[typed]{Supposons}}{Supposons}{insertion}{ega25_25_020-v0021} S=Spec(k), k un corps, \var{\inserted[typed]{et que}}{et que}{insertion}{ega25_25_020-v0022} \var{\xout{xxxxx}}{}{typewriter strike-out}{ega25_25_020-v0023} pour toute composante irréduc-%
tible  $X_i$ de X, on a\hand{it} dim $\hand{\overline{\textcolor{black}{f(X_i)}}} \hand{\geqslant} 2$, Alors, si X sur k satisfait $C_{m+1}$ ,
alors E\hand{$_{\xi}$} contient un ouvert dense U de P .

\mspage{\hand{20}}

% NOTE: A pen curve runs upward from the circled lower v/vi label to the preceding item.
